% Options for packages loaded elsewhere
\PassOptionsToPackage{unicode}{hyperref}
\PassOptionsToPackage{hyphens}{url}
\PassOptionsToPackage{dvipsnames,svgnames,x11names}{xcolor}
%
\documentclass[
  11pt,
]{article}
\usepackage{amsmath,amssymb}
\usepackage{iftex}
\ifPDFTeX
  \usepackage[T1]{fontenc}
  \usepackage[utf8]{inputenc}
  \usepackage{textcomp} % provide euro and other symbols
\else % if luatex or xetex
  \usepackage{unicode-math} % this also loads fontspec
  \defaultfontfeatures{Scale=MatchLowercase}
  \defaultfontfeatures[\rmfamily]{Ligatures=TeX,Scale=1}
\fi
\usepackage{lmodern}
\ifPDFTeX\else
  % xetex/luatex font selection
  \setmainfont[]{DejaVu Serif}
  \setmonofont[]{DejaVu Sans Mono}
\fi
% Use upquote if available, for straight quotes in verbatim environments
\IfFileExists{upquote.sty}{\usepackage{upquote}}{}
\IfFileExists{microtype.sty}{% use microtype if available
  \usepackage[]{microtype}
  \UseMicrotypeSet[protrusion]{basicmath} % disable protrusion for tt fonts
}{}
\makeatletter
\@ifundefined{KOMAClassName}{% if non-KOMA class
  \IfFileExists{parskip.sty}{%
    \usepackage{parskip}
  }{% else
    \setlength{\parindent}{0pt}
    \setlength{\parskip}{6pt plus 2pt minus 1pt}}
}{% if KOMA class
  \KOMAoptions{parskip=half}}
\makeatother
\usepackage{xcolor}
\usepackage[margin=1in]{geometry}
\setlength{\emergencystretch}{3em} % prevent overfull lines
\providecommand{\tightlist}{%
  \setlength{\itemsep}{0pt}\setlength{\parskip}{0pt}}
\setcounter{secnumdepth}{-\maxdimen} % remove section numbering
\ifLuaTeX
  \usepackage{selnolig}  % disable illegal ligatures
\fi
\usepackage{bookmark}
\IfFileExists{xurl.sty}{\usepackage{xurl}}{} % add URL line breaks if available
\urlstyle{same}
\hypersetup{
  pdftitle={Five-Colorings of Glued Triangular Lattices with Auxiliary Sites},
  pdfauthor={Aafi Khan Malek},
  colorlinks=true,
  linkcolor={Maroon},
  filecolor={Maroon},
  citecolor={Blue},
  urlcolor={Blue},
  pdfcreator={LaTeX via pandoc}}

\title{Five-Colorings of Glued Triangular Lattices with Auxiliary Sites}
\author{Aafi Khan Malek}
\date{November 2026}

\ifPDFTeX
\DeclareUnicodeCharacter{0393}{\ensuremath{\Gamma}}
\DeclareUnicodeCharacter{039B}{\ensuremath{\Lambda}}
\DeclareUnicodeCharacter{03B3}{\ensuremath{\gamma}}
\DeclareUnicodeCharacter{03B4}{\ensuremath{\delta}}
\DeclareUnicodeCharacter{03C1}{\ensuremath{\rho}}
\DeclareUnicodeCharacter{03C3}{\ensuremath{\sigma}}
\DeclareUnicodeCharacter{03C4}{\ensuremath{\tau}}
\DeclareUnicodeCharacter{03C9}{\ensuremath{\omega}}
\DeclareUnicodeCharacter{2084}{\ensuremath{{}_4}}
\DeclareUnicodeCharacter{2208}{\ensuremath{\in}}
\DeclareUnicodeCharacter{2209}{\ensuremath{\notin}}
\DeclareUnicodeCharacter{2212}{\ensuremath{-}}
\DeclareUnicodeCharacter{221A}{\ensuremath{\surd}}
\DeclareUnicodeCharacter{2229}{\ensuremath{\cap}}
\DeclareUnicodeCharacter{222A}{\ensuremath{\cup}}
\DeclareUnicodeCharacter{2260}{\ensuremath{\ne}}
\DeclareUnicodeCharacter{2282}{\ensuremath{\subset}}
\DeclareUnicodeCharacter{22C3}{\ensuremath{\bigcup}}
\DeclareUnicodeCharacter{2194}{\ensuremath{\leftrightarrow}}
\fi
\begin{document}
\maketitle

Independent Researcher, Ahmedabad, India ·
\href{mailto:aafimalek2032@gmail.com}{\nolinkurl{aafimalek2032@gmail.com}}.\\
This research note was developed with AI assistance.

This is a handwritten infinite structural obstruction for the
\href{https://github.com/Aafimalek/novel-discoveries/tree/a27775bee5782ed20476b0bec59d28be5031ccde/hadwiger-nelson/verification/parts-native-field-obstruction/exterior-cross-fixture-81}{81-site
fixture}, not for all Parts-type coordinate sets.

Let \texttt{s=√−3}, \texttt{ω=(1+s)/2}, \texttt{Λ=Z{[}ω{]}},
\texttt{K=Q(s,√33)}, and \texttt{J=Q(√−11)⊂K} where √−11=s√33/3. Let
\texttt{ρ=(7+√−15)/8}, the exterior rotation used by
\href{https://arxiv.org/html/2010.12665v2\#S4}{Parts,
arXiv:2010.12665v2, §§2,4}. In particular \textbar ρ\textbar=1 and
Re(ρ)=7/8.

For any set Γ of norm-one elements of J, put
\texttt{B(Γ)=\{0\}\ ∪\ ⋃\_\{γ∈Γ\}\ γΛ}, where γ and −γ represent the
same lattice. Allow arbitrary Γ, including infinite sets, and take
independent Γ\_A and Γ\_B. Define

\begin{verbatim}
U = B(Gamma_A)  union  rho B(Gamma_B).
\end{verbatim}

\textbf{Theorem.} The induced unit-distance graph on U admits a proper
five-coloring. For any chosen distance-two pair in U, it has a proper
five-coloring making those two colors equal. The latter assertion is
pointwise in the chosen pair, not a single palette that simultaneously
makes every such pair equal.

The 81-site fixture is contained in this set for \texttt{Γ\_A=\{1,σ,τ\}}
and \texttt{Γ\_B=\{1,σ\}}, where \texttt{σ=(5+√−11)/6} and
\texttt{τ=(49+3√−11)/50}. The theorem includes doubling all three full
lattices under ρ, or arbitrary further norm-one J orientations on either
side.

\subsection{Classification of cross
distances}\label{classification-of-cross-distances}

The square classes of -3,33,-15 are independent, using primes 11 and 5
and then the sign of -3. In particular j=√−15 is not in K and ρ∉K. Thus
the two sides of U share only zero.

For x∈B(Γ\_A) and y∈B(Γ\_B), expand the squared cross distance as

\begin{verbatim}
|x-rhoy|^2 = N(x)+N(y) - (7/4)Re(x*conj(y))
          + (j/8)(x*conj(y)-y*conj(x)),
\end{verbatim}

where N(x)=\textbar x\textbar². The terms besides the last lie in K and
the coefficient of j lies in K. Hence a rational squared distance forces
x·conj(y)=y·conj(x). If both x and y are nonzero, they are parallel, in
the sense of lying on the same real line through the origin, with either
sign allowed.

Parallel nonzero lattice sites must come from the same orientation up to
sign. To see this, suppose they come from γΛ and δΛ, so their quotient
has the form f g with f=γ/δ=a+b√−11∈J and g=c+d s∈Q(s) nonzero. If f is
not ±1, then b≠0 (a real norm-one element of J is ±1), and a≠0 (the
equation 11b²=1 has no rational b). The imaginary part of f g has
independent terms a d√3 and b c√11. It cannot vanish unless c=d=0. This
contradiction proves the claim.

For one orientation write x=γap, y=γbp with p∈Λ and a,b∈Z. Such a
lattice direction p exists because two real-parallel Eisenstein points
have rational quotient: the quotient belongs to Q(s) and its real
subfield is Q. One can take p primitive in the two integer coordinates,
so all points on that rational line are integral multiples of p.~The
exact norm identity is

\begin{verbatim}
16|x-rhoy|^2 = N(p) [(a+b)^2 + 15(a-b)^2],
\end{verbatim}

with positive integer N(p).

\begin{itemize}
\tightlist
\item
  At squared distance one the right side is 16. If a≠b, then
  \textbar a−b\textbar=1, N(p)=1 and \textbar a+b\textbar=1, forcing one
  endpoint zero. If a=b, then N(p)a²=4: the cross edge joins the same
  native point x=y on the radius-two shell. Edges involving zero already
  belong to one of the two sides.
\item
  At squared distance four the right side is 64. If a=b, N(p)a²=16 and
  both endpoint radii are four. If \textbar a−b\textbar=1,
  \textbar a+b\textbar{} must be one of 1,3,5,7; the corresponding
  values of (a+b)²+15 are 16,24,40,64. Only 16 and 64 divide 64. The
  first again has a zero endpoint; the last gives one radius three and
  one radius four, with N(p)=1. If \textbar a−b\textbar=2 then N(p)=1
  and \textbar a+b\textbar=2, again a zero endpoint. Larger differences
  are impossible.
\end{itemize}

Thus genuinely cross-side unit edges are exactly the matching radius-two
references where both matching points are present. A genuinely
cross-side distance-two pair has radii (4,4), (4,3), or (3,4).

\subsection{The shell is independent}\label{the-shell-is-independent}

In each triangular lattice γΛ, the radius-two shell is 2γ times the six
norm-one Eisenstein units. Indeed N(v)=4 implies v∈2Λ because Λ/2Λ is
the field F₄, conjugation descends to an automorphism of that field, and
the image of v·conj(v) is zero. Then v/2 has norm one. Solving
N(m+nω)=m²+mn+n²=1 gives the six units stated. Two distinct shell points
in the same orientation are at least two apart, because their directions
are multiples of 60 degrees.

For shell points in distinct orientations, a unit distance would require
Re(f w)=7/8, where f=a+b√−11 is their norm-one orientation ratio and w
is a sixth root of unity. If w is nonreal, its imaginary part is ±√3/2,
and this expression contains a nonzero irrational term ±b√33/2 in
addition to a rational number. If w=±1, it would require a=±7/8, hence
11b²=15/64. This has no rational solution: the right side has odd 5-adic
valuation whereas 11b² has even 5-adic valuation. Thus each side's
entire radius-two shell is independent.

\subsection{Explicit coloring, with optional terminal
equality}\label{explicit-coloring-with-optional-terminal-equality}

Use the
\href{https://github.com/Aafimalek/novel-discoveries/tree/b2f2d863aa4903b8aee52efc006e26f3c8b738b8/hadwiger-nelson/verification/parts-native-field-obstruction}{native
2-adic four-coloring} of the unit-distance graph on all of K, and copy
it onto the exterior side via inverse rotation. Both agree on zero.
Before recoloring, every unit edge inside each side is proper. Recolor
the entire radius-two shell on either one side with a fresh fifth color.
The shell is independent there, and every genuine cross-side unit edge
has exactly one endpoint recolored. This gives a proper five-coloring,
also for infinite orientation sets.

For a chosen distance-two pair contained in one side, recolor the
\emph{opposite} shell instead. In the side containing the pair the
native 2-adic coloring gives the two points equal colors; there is no
change at zero. For a genuinely cross-side distance-two pair, select a
radius-four endpoint and recolor that endpoint's side's shell. The
radius-four point is not adjacent to its own radius-two shell by the
reverse triangle inequality. The other endpoint is also not adjacent to
that shell: genuine cross-side unit edges have both endpoints at radius
two, whereas the other endpoint has radius three or four. Give the two
nonadjacent terminals the same fifth color as the shell. No edge becomes
monochromatic. This completes the proof.

\subsection{How to leave the
obstruction}\label{how-to-leave-the-obstruction}

The theorem covers unions of individual triangular lattice orientations,
not their Minkowski sums, additive closures, or arbitrary native-field
points. In particular, actual Parts-type auxiliary joins use different
nonzero collinear radii in opposite copies, which the theorem excludes
within the union family. The standard radii √11/2±√3/6 have irrational
ratio, while above any two parallel nonzero sites in the native
orientation union have rational ratio. Thus the auxiliary join is a
direct algebraic way to escape the original incidence classification. We
now prove that this auxiliary family still admits five colors.

This written result is not formalized in Lean.

\subsection{Extension to the standard unequal-radius auxiliary
sites}\label{extension-to-the-standard-unequal-radius-auxiliary-sites}

Write \(r=\sqrt{33}\), \(p=s(r+1)/6\), and \(q=s(r-1)/6\).

In each native side, start with any union of norm-one J rotations of
full Λ (possibly infinite and independent between sides). Also permit an
arbitrary set of auxiliary sites of the form γp or γq for any norm-one
γ∈J, without requiring that their orientation occur in the ordinary
lattice union. Allow signs independently (−γ is again norm one). Let A,B
be the resulting native sets, each including 0. Then the unit-distance
graph on

\begin{verbatim}
A  union  rhoB
\end{verbatim}

admits a proper five-coloring. For any selected pair at distance two
anywhere in this union, there is a proper five-coloring with that pair
equal. This is pointwise in the choice of pair.

In particular take the earlier infinite glued fixture and adjoin p,q to
each native side, i.e.~add p,q,ρp,ρq to the union. This gives genuine
unequal-radius cross-copy unit edges, but does not escape the coloring
obstruction.

\subsubsection{Exact cross-incidence
classification}\label{exact-cross-incidence-classification}

If x∈A and y∈B, the preceding expansion of \textbar x−ρy\textbar² shows
that any rational squared cross distance imposes x·conj(y)∈R, since
√−15∉K. Thus nonzero native representatives x,y must lie on the same
real line. The overlap between the two sides is still only zero.

\textbf{One auxiliary and one full-lattice point.} Cancel the auxiliary
orientation, which leaves a full-lattice point of the form f z with
\textbar f\textbar=1 in J and z∈Λ. Since p and q are purely imaginary, f
z must also be purely imaginary. Write f=a+b√−11, z=c+d s, with a,b,c,d
rational. If f is not ±1, both a and b are nonzero (there is no rational
b with 11b²=1), and

\begin{verbatim}
Re(f z) = ac - bdsqrt33.
\end{verbatim}

A zero real part then forces c=d=0 by irrationality of √33. For a
nonzero full-lattice point the orientation must therefore be f=±1. Also
Λ∩iR=Zs: in the lattice coordinates m+nω, zero real part means 2m+n=0.
Thus the possible native pair, after dividing out the common rotation,
is ks and s(r±1)/6 for an integer k. Either order of the two points
gives the same squared cross norm. It equals

\begin{verbatim}
3 [ k^2 + ((r+/-1)/6)^2 - (7/4)k(r+/-1)/6 ].
\end{verbatim}

The coefficient of r is ±1/6−7k/8, which cannot vanish for integer k.
This also holds at k=0. Hence no such cross pair has even a rational
squared distance. In particular there are no unit edges or distance-two
pairs of this type.

\textbf{Two auxiliary representatives.} Their rays can be parallel only
if their relative orientation in J is real (the ratio p/q is real),
hence ±1. Fix one common orientation and keep the possible signs
explicitly. Then N(p)=(17+r)/6, N(q)=(17−r)/6, N(p)+N(q)=17/3 and
p·conj(q)=8/3. Matching radial types p/p or q/q have squared cross norm
equal to N(p) or N(q) times 1/4 (same sign) or 15/4 (opposite signs);
these are irrational. Mixed p/q with opposite signs has squared norm
17/3+14/3=31/3. Mixed p/q with the same sign has squared norm
17/3−14/3=1. As a consequence the new cross edges are exactly

\begin{verbatim}
gammap  <->  rhogammaq   and   gammaq  <->  rhogammap
\end{verbatim}

when the endpoints are present, with γ including the common sign. There
are no new cross distance-two pairs.

\textbf{Two full-lattice representatives.} The earlier classification
continues to apply: away from shared origin, cross unit edges are
matching radius-two references, with matching native representatives,
and any genuinely cross-side distance-two pair has radii (4,4), (4,3) or
(3,4).

\subsubsection{Coloring}\label{coloring}

Give each side the same
\href{https://github.com/Aafimalek/novel-discoveries/tree/b2f2d863aa4903b8aee52efc006e26f3c8b738b8/hadwiger-nelson/verification/parts-native-field-obstruction}{native
2-adic coloring} after applying inverse rotation on the exterior side.
It is already proper on every edge internal to a side, including any
auxiliary vertex. Moreover, it distinguishes each newly introduced mixed
auxiliary edge across sides: the difference of the two native
representatives is ±γ(p−q)=±γs/3 of squared norm 1/3. Conjugation
preserves the normalized valuation used by the native coloring; this
norm has valuation zero, so the difference itself has valuation zero. It
lies in the integral coset and its F₄ residue is nonzero, exactly the
certificate for different native colors. The two sides still agree at 0.

Neither auxiliary radius is two, so both radius-two shells are the
independent full-lattice shells proved earlier. Recolor either of those
shells with a fresh fifth color. This fixes the matching reference edges
without disturbing any new auxiliary cross edge, proving proper
five-colorability.

For two chosen distance-two points contained in one side, recolor the
other shell instead. The native coloring gives equal colors to such a
pair in K; this also covers pairs involving the shared 0. For a
genuinely cross-side pair at distance two, the classification is
unchanged. Recolor the shell on a radius-four endpoint's side; this
point cannot neighbor its own shell (radial gap two), and the other
terminal cannot neighbor that shell across sides (only the reference and
auxiliary unit incidences above exist). Give both terminals the fresh
fifth color. They are not neighbors of each other, so this is a valid
equality witness.

This result uses only unions of single rotated lattices plus the stated
auxiliary radius types. It makes no claim about adding arbitrary native
points or Minkowski sums, or about the published full five-chromatic
Parts graphs. These are written proofs; they have not been formalized in
Lean.

\end{document}
